\section{Monolith y Blast Radius: el límite del código no es el límite de las fallas}

Asif no considera que los microservice sean un requisito para escalar. La clase se centra más en el \term{blast radius} (radio de explosión): a cuántos usuarios, funciones y region puede afectar una falla de código, de configuración o de una dependencia, y durante cuánto tiempo. Un monolith puede compartir el código y el proceso de despliegue y, al mismo tiempo, aislar las rutas de alto riesgo mediante traffic cell, queue, database, feature flag y límites de permisos.

\subsection{Separar la critical path de la experimental path}

Esta sección responde a la pregunta «cómo puede un solo repositorio seguir teniendo dominios de falla pequeños». Al leer la figura, no cuente el número de service; revise los límites de ejecución de auth, editor request, index job, new model y shared schema. Si el experimental code puede agotar las conexiones de la critical database o contaminar una queue compartida, el blast radius sigue siendo grande aunque ese código se haya separado en un service independiente.

\lecturefigure{03-blast-radius.png}{Un monolith puede reducir el blast radius mediante cells de despliegue y límites de dependencias y de tráfico.}{Subtítulos locales 05:20--07:50; redibujo conceptual.}

\begin{knowledgebox}{Lectura de la figura: cuatro controles de aislamiento}
\textbf{Traffic} se aísla con canary y cells por tenant/region; \textbf{resource}, con quota, queue y conjunto independientes; \textbf{state}, con versioned schema y rebuildable derived data; \textbf{change}, con feature flag, rollback y deploy policy. Los microservice son solo una de las formas de implementar estos aislamientos.
\end{knowledgebox}

\begin{importantbox}{Fórmula de diagnóstico del blast radius}
Puede escribirse de forma aproximada como
\[
B=A\times D\times S,
\]
donde $A$ son las affected requests/users, $D$ es la duración y $S$ es el daño al estado o la complejidad de la recuperación. La arquitectura debe reducir primero la $S$ irreversible; después, reducir $A$ con canary y cells, y reducir $D$ con detection/rollback.
\end{importantbox}

\begin{warningbox}{«Mantener el código simple» no significa escribir menos abstracciones}
La simplicidad debe reflejarse en invariant explicables, pocas dependency, failure mode que puedan probarse y recovery path claras. Si la complejidad se esconde en global state, en retry implícitos o en una database compartida, el número de líneas de código puede disminuir, pero la complejidad del sistema aumenta.
\end{warningbox}

\teachervoice{La postura de Asif no es «nunca separar servicios», sino dar prioridad a que un grupo reducido de ingenieros entienda la critical path completa. Solo cuando la organización, los dominios de falla o la forma de los recursos realmente necesiten evolucionar de manera independiente conviene pagar el costo de los límites distribuidos.}

\subsection{Resumen de la sección}

Escalar es, ante todo, un problema de aislamiento, no de número de service. Tanto un buen monolith como un buen distributed system necesitan un blast radius explícito, presupuestos de recursos, contratos de versión y estrategias de recuperación.
